
\documentclass{article}
\usepackage{colortbl}
\usepackage{makecell}
\usepackage{multirow}
\usepackage{supertabular}

\begin{document}

\newcounter{utterance}

\centering \large Interaction Transcript for game `hot\_air\_balloon', experiment `air\_balloon\_survival\_hr\_complexity\_hard', episode 0 with mistral{-}large{-}3{-}azure.
\vspace{24pt}

{ \footnotesize  \setcounter{utterance}{1}
\setlength{\tabcolsep}{0pt}
\begin{supertabular}{c@{$\;$}|p{.15\linewidth}@{}p{.15\linewidth}p{.15\linewidth}p{.15\linewidth}p{.15\linewidth}p{.15\linewidth}}
    \# & $\;$A & \multicolumn{4}{c}{Game Master} & $\;\:$B\\
    \hline

    \theutterance \stepcounter{utterance}  
    & & \multicolumn{4}{p{0.6\linewidth}}{
        \cellcolor[rgb]{0.9,0.9,0.9}{
            \makecell[{{p{\linewidth}}}]{
                \texttt{\tiny{[P1$\langle$GM]}}
                \texttt{Sudjelujete u suradničkoj pregovaračkoj igri.} \\
\\ 
\texttt{Zajedno s drugim sudionikom morate se dogovoriti o jednom skupu predmeta koji će se zadržati. Svaki od vas ima vlastiti pogled na to koliko vam je svaki predmet važan (važnost). Ne znate kako drugi sudionik vrednuje predmete. Dodatno, dobivate i napor koji svaki predmet zahtijeva.} \\
\texttt{Možete se dogovoriti o skupu samo ako ukupni napor odabranih predmeta ne premašuje zajedničko ograničenje:} \\
\\ 
\texttt{OGRANIČENJE = 4734} \\
\\ 
\texttt{Ovdje su pojedinačne vrijednosti napora za predmete:} \\
\\ 
\texttt{Napor predmeta = \{"C08": 192, "C55": 325, "A43": 812, "A67": 868, "B42": 593, "A62": 915, "B19": 944, "C43": 311, "C73": 252, "C10": 343, "B36": 104, "A18": 558, "C61": 627, "C67": 593, "C42": 827, "C70": 611, "B40": 95, "C39": 251, "A21": 226, "C16": 21\}} \\
\\ 
\texttt{Ovdje je vaš osobni pogled na važnost svakog predmeta:} \\
\\ 
\texttt{Vrijednosti važnosti predmeta = \{"C08": 286, "C55": 419, "A43": 906, "A67": 962, "B42": 687, "A62": 1009, "B19": 1038, "C43": 405, "C73": 346, "C10": 437, "B36": 198, "A18": 652, "C61": 721, "C67": 687, "C42": 921, "C70": 705, "B40": 189, "C39": 345, "A21": 320, "C16": 115\}} \\
\\ 
\texttt{Cilj:} \\
\\ 
\texttt{Vaš je cilj pregovarati o zajedničkom skupu predmeta koji vam koristi što je više moguće (tj. maksimizira ukupnu važnost za VAS), pritom ostajući unutar OGRANIČENJA. Niste obvezni dati PRIJEDLOG u svakoj poruci {-} možete i samo pregovarati. Sve taktike su dopuštene!} \\
\\ 
\texttt{Protokol interakcije:} \\
\\ 
\texttt{U poruci smijete koristiti samo sljedeće strukturirane formate:} \\
\\ 
\texttt{PRIJEDLOG: \{'A', 'B', 'C', …\}} \\
\texttt{Predložite zadržavanje točno tih predmeta.} \\
\texttt{ODBIJ: \{'A', 'B', 'C', …\}} \\
\texttt{Izričito odbijte protivnikov prijedlog.} \\
\texttt{ARGUMENT: \{'...'\}} \\
\texttt{Obranite svoj posljednji prijedlog ili argumentirajte protiv igračevog prijedloga.} \\
\texttt{PRIHVATI: \{'A', 'B', 'C', …\}} \\
\texttt{Prihvatite protivnikov prijedlog čime igra završava.} \\
\\ 
\\ 
\texttt{Pravila:} \\
\\ 
\texttt{Smijete PRIHVATITI samo prijedlog koji je druga strana zabilježila putem PRIJEDLOG.} \\
\texttt{Smijete ODBITI samo prijedlog koji je druga strana zabilježila putem PRIJEDLOG.} \\
\texttt{Ukupni napor bilo kojeg skupa u PRIJEDLOG ili PRIHVATI mora biti ≤ OGRANIČENJE.} \\
\texttt{NE otkrivajte svoje skrivene ocjene važnosti.} \\
\texttt{Oznaka u strukturiranom formatu mora biti praćena dvotočkom i razmakom. Argument mora biti python skup koji sadrži 0 ili više nizova znakova.} \\
\texttt{Dakle, mora biti u obliku "OZNAKA: \{...\}"} \\
\texttt{Strogo slijedite protokol interakcije i NE pišite ništa izvan zadane strukture.} \\
\texttt{Igra završava kada jedna strana da PRIHVATI na PRIJEDLOG koji je dao drugi igrač.} \\
\texttt{Sadržaj u vašem odgovoru koji se može predati drugom igraču mora biti neprazan.} \\
\texttt{Aktivni su samo prijedlozi koji su zabilježeni putem strukturiranog formata PRIJEDLOG i koji nisu odbijeni putem ODBIJ.} \\
\\ 
            }
        }
    }
    & \\ \\

    \theutterance \stepcounter{utterance}  
    & \multicolumn{4}{p{0.6\linewidth}}{
        \cellcolor[rgb]{0.9,0.9,0.9}{
            \makecell[{{p{\linewidth}}}]{
                \texttt{\tiny{[P1$\rangle$GM]}}
                \texttt{Kako bih maksimizirao svoju ukupnu važnost unutar zadanog ograničenja napora (4734), predložit ću skup predmeta koji su za mene visoko vrijedni, a imaju relativno nizak napor. Evo mog početnog prijedloga:} \\
\\ 
\texttt{**PRIJEDLOG: \{'C08', 'C55', 'C73', 'C10', 'B36', 'A18', 'C61', 'C67', 'B40', 'C39', 'A21', 'C16'\}**} \\
\\ 
\texttt{Objašnjenje (koje mogu koristiti u pregovorima ako bude potrebno):} \\
\texttt{{-} Skup je pažljivo odabran da ostane unutar ograničenja (ukupni napor = 4728, što je ≤ 4734).} \\
\texttt{{-} Fokusiran je na predmete s visokim omjerom važnosti i napora (npr. *C16* ima vrlo nizak napor, a *A18* i *C61* su visoko vrednovani).} \\
\texttt{{-} Izbjegnuti su predmeti s visokim naporom osim ako nisu izuzetno važni (npr. *A62* i *B19* su isključeni jer imaju visok napor).} \\
\\ 
\texttt{Ako druga strana ima primjedbe, otvoren sam za pregovore – možemo razmjenjivati predmete dok ne postignemo kompromis.} \\
            }
        }
    }
    & & \\ \\

    \theutterance \stepcounter{utterance}  
    & & \multicolumn{4}{p{0.6\linewidth}}{
        \cellcolor[rgb]{0.9,0.9,0.9}{
            \makecell[{{p{\linewidth}}}]{
                \texttt{\tiny{[P1$\langle$GM]}}
                \texttt{Vaš odgovor sadržavao je neoznačeni niz ili ste upotrijebili STRATEŠKO RAZMIŠLJANJE više od jednom. Smijete koristiti samo strukturirane formate kako je objašnjeno u početnoj poruci. Svi moraju biti u obliku 'TAG: \{...\}'.} \\
            }
        }
    }
    & \\ \\

    \theutterance \stepcounter{utterance}  
    & \multicolumn{4}{p{0.6\linewidth}}{
        \cellcolor[rgb]{0.9,0.9,0.9}{
            \makecell[{{p{\linewidth}}}]{
                \texttt{\tiny{[P1$\rangle$GM]}}
                \texttt{**PRIJEDLOG: \{'C08', 'C55', 'A43', 'A67', 'B42', 'A18', 'C61', 'C67', 'C10', 'B36', 'C39', 'A21', 'C16'\}**} \\
\\ 
\texttt{Ukupni napor: **4727** (unutar ograničenja).} \\
            }
        }
    }
    & & \\ \\

    \theutterance \stepcounter{utterance}  
    & & & \multicolumn{2}{p{0.3\linewidth}}{
        \cellcolor[rgb]{0.9,0.9,0.9}{
            \makecell[{{p{\linewidth}}}]{
                \texttt{\tiny{[GM$|$GM]}}
                \texttt{abort} \\
            }
        }
    }
    & & \\ \\

    \theutterance \stepcounter{utterance}  
    & & & \multicolumn{2}{p{0.3\linewidth}}{
        \cellcolor[rgb]{0.9,0.9,0.9}{
            \makecell[{{p{\linewidth}}}]{
                \texttt{\tiny{[GM$|$GM]}}
                \texttt{end game} \\
            }
        }
    }
    & & \\ \\

\end{supertabular}
}

\end{document}
